%---------------------------------------------------------------------------%
%->> Frontmatter
%---------------------------------------------------------------------------%
%-
%-> 封面与声明
%-
% 说明：非个人成果前置研究稿，采用中性标题页，不渲染个人身份与原创声明页
\maketitle% 生成标题页
%\MAKETITLE% 英文标题页
%\makedeclaration% 不渲染个人原创声明页

%-
%-> 中文摘要
%-
\intobmk\chapter*{摘\quad 要}% 显示在书签但不显示在目录
\setcounter{page}{1}% 开始页码
\pagenumbering{Roman}% 页码符号

本稿为低能原子核结构中原子核形状与集体激发的综合前置研究稿。低能核物理以核子—核子有效相互作用与微观多体关联为出发点，研究有限量子多体系统在实验室系中表现出的转动、振动与空间几何形状效应。本文围绕“原子核几何形状在量子力学意义下的界定、微观产生机制、典型质量区模式竞争及其在束谱学实验验证”这一主线展开。全文设计为五章：从实验能谱的多重解释与形状概念的物理本质出发，系统建立微观壳模型、平均场形变、配对准粒子与对称性恢复理论语言；深入讨论三轴形变下的摆动运动、核手征性及多模式耦合动力学；在六个代表性质量区开展唯象比较与探针约束分析；系统梳理重离子熔合蒸发、数字成形脉冲处理、在束伽马符合谱学与跃迁几率测量实验技术链。

本稿包含五章、三十九节及六篇配套附录，作为博士论文前置非个人成果部分的研究稿。正文区分实验报告、模型解释与条件化推演；附录给出符号约定、推导、数值控制及图表依据。个人实验、装置参数、数据与成果须依据实际研究另行补充，本稿不作相关归属或学校送审状态的声明。

\keywords{原子核形状，三轴形变，集体转动，核手征性，在束伽马谱学}% 中文关键词
%-
%-> 英文摘要
%-
\pagestyle{frontmatterstyle}%
\intobmk\chapter*{Abstract}% 显示在书签但不显示在目录
\pagestyle{enfrontmatterstyle}%

This document serves as a pre-study research draft investigating nuclear shapes and collective excitations within low-energy nuclear structure physics. Rooted in effective nucleon-nucleon interactions and microscopic many-body correlations, low-energy nuclear physics investigates how rotational, vibrational, and spatial geometric phenomena emerge in finite quantum systems within the laboratory frame. This work systematically examines the quantum definition of nuclear shapes, microscopic driving mechanisms, collective mode competition across representative mass regions, and experimental in-beam spectroscopy validation.

The overall architecture spans five chapters: clarifying the relationship between discrete gamma-ray spectra and underlying spatial configurations, formulating mean-field deformation, pairing quasiparticles, and symmetry restoration; analyzing triaxial wobbling, nuclear chirality, and mode coupling dynamics; comparing empirical signatures across six characteristic mass regions; and detailing the full experimental measurement chain from heavy-ion fusion-evaporation reactions to digital pulse processing, coincidence matrices, and absolute transition probabilities. The manuscript contains five chapters, thirty-nine sections and six supporting appendices. It distinguishes experimental reports, model interpretations and conditional derivations, with conventions, numerical controls and figure provenance documented separately. Personal experiments, apparatus parameters, data and results require the corresponding actual research; this draft makes no claim to those results or to readiness for institutional examination.

\KEYWORDS{nuclear shape, triaxial deformation, collective rotation, nuclear chirality, in-beam gamma spectroscopy}% 英文关键词

\cleardoublepage\pagestyle{frontmatterstyle}%
%---------------------------------------------------------------------------%
